\documentclass[11pt]{article}

\usepackage[T1]{fontenc}
\usepackage{lmodern}
\usepackage[margin=1.08in]{geometry}
\usepackage{amsmath,amssymb,amsthm,mathtools}
\usepackage{xcolor}
\usepackage{microtype}
\usepackage[colorlinks=true,linkcolor=blue!55!black,
            citecolor=blue!55!black,urlcolor=blue!65!black]{hyperref}

\newtheorem{theorem}{Theorem}[section]
\newtheorem{proposition}[theorem]{Proposition}
\newtheorem{lemma}[theorem]{Lemma}
\newtheorem{corollary}[theorem]{Corollary}
\theoremstyle{remark}
\newtheorem{remark}[theorem]{Remark}

\newcommand{\R}{\mathbb R}
\newcommand{\ord}{\operatorname{ord}}
\newcommand{\Hess}{\operatorname{Hess}}
\newcommand{\cO}{\mathcal O}

\title{All-order compatibility of the resonant Taylor equation at an
umbilic of a \(W\)-congruence}
\author{[Author Name]}
\date{August 18, 2026}

\begin{document}
\maketitle

\begin{abstract}
\textbf{Relation to the problem list.}
This manuscript addresses Question~8 (L8), ``All-order identity for umbilics
of $W$-congruences.''  It gives a full proof of the nontrivial finite-jet
interpretation: the resonant equation follows from the lower-order jet
equations.  For an already existing analytic $W$-congruence the literal
identity is automatic from $W\equiv0$; the result proved here is precisely the
claimed redundancy in the Taylor recursion.

Craizer and Garcia isolated one apparently exceptional equation in the
Taylor recursion at a nondegenerate umbilic of a \(W\)-congruence.  If the
integer determined by the quadratic jet is \(m\), their conjecture says that
the resonant equation \(W_{m0}=0\) imposes no condition beyond the equations
of lower total order.  We prove this in every order.  In hodograph variables
the Weingarten equation is equivalent, up to a nonvanishing factor, to a
function \(E\), and there is an exact identity
\[
                         (X-2)\Delta=A E.
\]
The linear part of \(X\) has eigenvalues \(2/(m+1)\) and \(1\).  A finite
Poincar\'e--Dulac reduction has at most one nonlinear resonance.  When that
resonance occurs, the identity above and one lower-degree equation force its
coefficient to vanish.  The remaining scalar resonance then forces the
coefficient of \(A^m\) in \(E\), and hence \(W_{m0}\), to vanish.  The
argument is finite and algebraic; no convergence of a formal normalizing
map is required.
\end{abstract}

\section{Statement of the jet problem}

Write \(P=\xi^1\) and \(Q=\xi^2\).  In the affine chart of
Craizer--Garcia, the equation for a \(W\)-congruence is
\begin{align}
 W(P,Q)={}&(Q_y-P_x)(Q_{xx}P_{yy}-P_{xx}Q_{yy}) \notag\\
 &-2P_y(P_{xx}Q_{xy}-Q_{xx}P_{xy})
 +2Q_x(P_{xy}Q_{yy}-P_{yy}Q_{xy}).                 \label{eq:W}
\end{align}
For a germ at the origin, set
\[
 p_{ij}=\partial_x^i\partial_y^jP(0),\qquad
 q_{ij}=\partial_x^i\partial_y^jQ(0),\qquad
 W_{ij}=\partial_x^i\partial_y^jW(P,Q)(0).
\]
Throughout, \(\cO(k)\) denotes terms of ordinary total order at least \(k\)
at the origin.
At an umbilic,
\[
 p_{01}=q_{10}=0,\qquad p_{10}=q_{01}.
\]
The common value in the second equality may be set equal to zero without
changing \(W\).  At a nondegenerate umbilic, the linear normalizations in
\cite{CG} give
\begin{equation}
 p_{20}=p_{02}=q_{11}=0,\qquad p_{11}=a\ne0,
 \qquad q_{02}\ne0,\qquad p_{11}\ne q_{02}.        \label{eq:norm}
\end{equation}
The parameter used in the source is
\begin{equation}
                              m=-\frac{q_{02}}a.       \label{eq:m}
\end{equation}

The literal assertion that \(W_{m0}=0\) for an already existing analytic
\(W\)-congruence follows at once from \(W\equiv0\).  The nontrivial content
of the conjecture is the following finite-jet implication.

\begin{theorem}[All-order resonant compatibility]                 \label{thm:main}
Assume \eqref{eq:norm}, and suppose that the parameter in
\eqref{eq:m} is an integer \(m\ge1\).  If
\begin{equation}
                  W_{ij}=0\qquad\text{whenever }i+j<m,             \label{eq:lower}
\end{equation}
then
\begin{equation}
                              W_{m0}=0.                             \label{eq:target}
\end{equation}
Thus the exceptional equation is a consequence of the lower total-order
equations and imposes no additional vanishing condition on the jet.
\end{theorem}

The proof occupies Sections~\ref{sec:hodograph}--\ref{sec:proof}.  Only
jets of finite order enter, so the theorem holds equally for formal germs
and for \(C^{m+2}\) germs.

For later reference, expanding a determinant gives another concise form of
\eqref{eq:W}:
\begin{equation}
 W=\det\!\begin{pmatrix}
 2Q_x&Q_y-P_x&-2P_y\\
 P_{xx}&P_{xy}&P_{yy}\\
 Q_{xx}&Q_{xy}&Q_{yy}
 \end{pmatrix}
 =\det\bigl(\Hess(xQ-yP),\Hess P,\Hess Q\bigr).       \label{eq:hess}
\end{equation}
In the last determinant, a symmetric Hessian is represented by the row
\((f_{xx},f_{xy},f_{yy})\).  Indeed,
\[
 \Hess(xQ-yP)=x\Hess Q-y\Hess P
                 +(2Q_x,Q_y-P_x,-2P_y),
\]
and row operations prove the second equality in \eqref{eq:hess}.

\section{Hodograph form of the equation}                         \label{sec:hodograph}

Assume for the moment that \(m>1\).  The first lower jet equation gives
\begin{equation}
                       W_{10}=2q_{20}a^2(1-m),                     \label{eq:W10}
\end{equation}
so \eqref{eq:lower} implies \(q_{20}=0\).

Introduce
\begin{equation}
 A=P_y,\qquad B=\frac{P_x-Q_y}{2}.                                 \label{eq:AB}
\end{equation}
At the origin,
\begin{equation}
 A=ax+\cO(2),\qquad
 B=\frac{(m+1)a}{2}\,y+\cO(2).                                   \label{eq:ABlinear}
\end{equation}
Consequently \((A,B)\) are local coordinates.  Denote their Jacobian by
\begin{equation}
 J=\det\frac{\partial(A,B)}{\partial(x,y)},\qquad
 J(0)=\frac{m+1}{2}a^2\ne0.                                      \label{eq:J}
\end{equation}
Regard the remaining two combinations of first derivatives as functions of
\((A,B)\):
\begin{equation}
 -Q_x=F(A,B),\qquad \frac{P_x+Q_y}{2}=G(A,B).                      \label{eq:FG}
\end{equation}
Equivalently,
\begin{equation}
 P_x=B+G,\qquad P_y=A,\qquad Q_x=-F,\qquad Q_y=G-B.                \label{eq:first}
\end{equation}

\begin{lemma}[Exact hodograph identities]                         \label{lem:identities}
Define
\begin{align}
 E={}&(1-G_B)F-(1+G_B)AF_A+(AG_A-B)F_B-2BG_A,          \label{eq:E}\\
 X={}&A(1+G_B)\partial_A+(B-AG_A)\partial_B,           \label{eq:X}\\
 \Delta={}&B^2-AF.                                    \label{eq:Delta}
\end{align}
Then
\begin{equation}
                         W=2J E                                      \label{eq:WE}
\end{equation}
and
\begin{equation}
                         (X-2)\Delta=A E.                             \label{eq:key}
\end{equation}
\end{lemma}

\begin{proof}
The following calculation makes the signs and the factor \(2\) explicit.
From \eqref{eq:first}, equality of mixed derivatives gives
\[
                         A_x-B_y=G_y,
        \qquad          F_y-B_x=-G_x.                              \label{eq:compat}
\]
A direct rearrangement of the six terms in \eqref{eq:W} yields
\begin{align}
 \frac W2={}&(A_x-B_y)\Delta_x+(F_y-B_x)\Delta_y \notag\\
 &+\det\!\begin{pmatrix}
 A&B&F\\ A_x&B_x&F_x\\ A_y&B_y&F_y
 \end{pmatrix}.                                                    \label{eq:Wdelta}
\end{align}
With the Poisson bracket
\(\{U,V\}=U_AV_B-U_BV_A\), the first line on the right of
\eqref{eq:Wdelta} is \(-J\{G,\Delta\}\).  The determinant in its second
line equals
\[
                         J(F-AF_A-BF_B).
\]
Consequently
\[
 \frac{W}{2J}=F-AF_A-BF_B-\{G,B^2-AF\}
\]
and expanding the bracket gives
\[
 \frac{W}{2J}=(1-G_B)F-(1+G_B)AF_A
                    +(AG_A-B)F_B-2BG_A,
\]
which is \eqref{eq:WE}.  This computation is valid near the origin because
\(J\) is a unit.

For the second identity, direct differentiation gives
\begin{align*}
 X(B^2-AF)-2(B^2-AF)
  =A\bigl((1-G_B)F-(1+G_B)AF_A
                    +(AG_A-B)F_B-2BG_A\bigr),
\end{align*}
which is \eqref{eq:key}.
\end{proof}

The normalized quadratic jet determines the leading terms of the objects in
Lemma~\ref{lem:identities}.  Since \(q_{20}=q_{11}=0\), the function \(F\)
has no constant or linear part.  From \eqref{eq:ABlinear} and
\eqref{eq:FG},
\begin{equation}
 F=\cO(2),\qquad
 G=\frac{1-m}{1+m}B+\cO(2).                           \label{eq:leadingFG}
\end{equation}
It follows that
\begin{equation}
 X=X_0+\cO(2),\qquad
 X_0=\alpha A\partial_A+B\partial_B,\qquad
 \alpha=\frac{2}{m+1},\qquad
 \Delta=B^2+\cO(3).                                  \label{eq:leading}
\end{equation}

\section{A finite resonance lemma}                                  \label{sec:normal}

We isolate the formal-algebraic step.  If \(H\) is a germ, \(H_d\) denotes
its homogeneous part of ordinary total degree \(d\).

\begin{lemma}[Vanishing of the resonant forcing]                    \label{lem:resonance}
Let \(m\ge2\), put \(\alpha=2/(m+1)\), and let
\[
 X=\alpha A\partial_A+B\partial_B+\cO(2),\qquad
 \Delta=B^2+\cO(3).
\]
If
\begin{equation}
                             R:=(X-2)\Delta=\cO(m+1),                \label{eq:Rorder}
\end{equation}
then the coefficient of \(A^{m+1}\) in \(R_{m+1}\) is zero.
\end{lemma}

\begin{proof}
We use only changes of coordinates tangent to the identity and only through
degree \(m\).  More explicitly, under old variables
\(z=\Phi(\widetilde z)\) we replace \(X\) by \(\Phi^*X\), \(\Delta\) by
\(\Delta\circ\Phi\), and \(R\) by \(R\circ\Phi\).  Thus the scalar identity
\((X-2)\Delta=R\) is preserved.  Since \(\Phi=\mathrm{id}+\cO(2)\) and
\(R=\cO(m+1)\), these changes preserve \(\Delta_2=B^2\), preserve the order
condition \eqref{eq:Rorder}, and leave the degree-\((m+1)\) homogeneous
part \(R_{m+1}\) unchanged.

Write \(X=X_0+\sum_{d\ge2}X_d\), where
\(X_0=\alpha A\partial_A+B\partial_B\) and the coefficients of \(X_d\)
are homogeneous of degree \(d\).  At degree \(d\), a coordinate change
generated by a homogeneous vector field \(H_d\) changes \(X_d\) by a signed
homological term.  With the pullback convention above it is
\(X_d\mapsto X_d-[X_0,H_d]\).  The sign is immaterial for elimination.  On
the monomial vector fields the operator \([X_0,\cdot]\) is diagonal:
\begin{align}
 [X_0,A^iB^j\partial_A]
    &=\bigl(\alpha(i-1)+j\bigr)A^iB^j\partial_A,       \label{eq:homA}\\
 [X_0,A^iB^j\partial_B]
    &=\bigl(\alpha i+j-1\bigr)A^iB^j\partial_B.       \label{eq:homB}
\end{align}
Here \(i+j=d\ge2\).  The eigenvalue in \eqref{eq:homA} never vanishes.
The eigenvalue in \eqref{eq:homB} can vanish only if
\[
             j=0,\qquad i=r:=\frac{m+1}{2}.
\]
Thus there is no nonlinear vector-field resonance when \(m\) is even.  If
\(m\) is odd, the only possible one is
\(cA^r\partial_B\), at degree \(r\).

Suppose \(m\) is odd.  Normalize the terms of \(X\) through degree \(r\),
leaving only that possible resonant term.  Since \(r+1<m+1\), the
degree-\((r+1)\) part of \eqref{eq:Rorder} is zero.  Using
\(\Delta_2=B^2\), it is
\begin{equation}
 0=(X_0-2)\Delta_{r+1}+cA^r\partial_B(B^2)
   =(X_0-2)\Delta_{r+1}+2cA^rB.                       \label{eq:killc}
\end{equation}
But
\[
                         \alpha r+1=2,
\]
so the coefficient of \(A^rB\) in
\((X_0-2)\Delta_{r+1}\) is zero.  Taking that coefficient in
\eqref{eq:killc} gives \(2c=0\).  Hence the sole possible nonlinear
resonance is absent.

We may therefore continue the finite normalization and arrange
\begin{equation}
                              X=X_0+\cO(m+1).                         \label{eq:Xnf}
\end{equation}
No convergence question arises: only finitely many homogeneous coordinate
changes have been used.  Since \(\Delta\) has order two, the remainder in
\eqref{eq:Xnf} contributes first in degree \(m+2\).  Consequently
\begin{equation}
                         R_{m+1}=(X_0-2)\Delta_{m+1}.                 \label{eq:Rtop}
\end{equation}
Finally, \(X_0-2\) acts on \(A^{m+1}\) by multiplication by
\[
                         \alpha(m+1)-2=0.
\]
The coefficient of \(A^{m+1}\) on the right-hand side of
\eqref{eq:Rtop}, and therefore in \(R_{m+1}\), is zero.
\end{proof}

\begin{remark}
The odd-\(m\) step \eqref{eq:killc} is essential.  A bare
Poincar\'e--Dulac argument would leave the term \(A^{(m+1)/2}\partial_B\).
Identity \eqref{eq:key} shows that its coefficient is already killed by an
equation of degree \((m+3)/2<m+1\).  This is the lower mixed compatibility
that is hidden in a direct Taylor elimination.
\end{remark}

\section{Proof of the conjecture}                                  \label{sec:proof}

\begin{proof}[Proof of Theorem~\ref{thm:main}]
First let \(m>1\).  Equation \eqref{eq:W10} and the lower-order hypothesis
give \(q_{20}=0\), so all conclusions of Section~\ref{sec:hodograph} apply.

Hypothesis \eqref{eq:lower} is precisely
\begin{equation}
                              W=\cO(m)                               \label{eq:Word}
\end{equation}
in the original variables.  The coordinate map \((x,y)\mapsto(A,B)\) is a
local diffeomorphism and \(J(0)\ne0\).  Hence \eqref{eq:WE} implies
\begin{equation}
                              E=\cO(m).                               \label{eq:Eord}
\end{equation}
By \eqref{eq:key},
\[
                           R:=(X-2)\Delta=AE=\cO(m+1).
\]
Lemma~\ref{lem:resonance}, together with \eqref{eq:leading}, now gives
\begin{equation}
                              [A^{m+1}](AE)_{m+1}=0,                  \label{eq:AEm}
\end{equation}
where \([A^{m+1}]\) denotes coefficient extraction.  Since multiplication
by \(A\) raises total degree by one and \(E=\cO(m)\), we have
\((AE)_{m+1}=AE_m\) and
\([A^{m+1}](AE)_{m+1}=[A^m]E_m\).  Thus \eqref{eq:AEm} says
\begin{equation}
                              [A^m]E_m=0.                             \label{eq:Em}
\end{equation}

It remains only to translate this coefficient back to the original jet.
First express \(W\) in the hodograph variables \((A,B)\), and denote its
degree-\(m\) part there by \(\widetilde W_m\).  Because \(E=\cO(m)\),
equation \eqref{eq:WE} gives
\(\widetilde W_m=2J(0)E_m\).  Pull this homogeneous polynomial back by the diagonal
linear map from \eqref{eq:ABlinear},
\[
                   A=ax,\qquad B=\frac{m+1}{2}ay.
\]
A monomial containing \(B\) cannot contribute to the coefficient of
\(x^m\).  With ordinary monomial coefficients, the exact relation is
\begin{equation}
 W_{m0}=2J(0)m!a^m[A^m]E_m.                            \label{eq:return}
\end{equation}
Equation \eqref{eq:Em} therefore gives \(W_{m0}=0\).

For \(m=1\), no hodograph normalization is needed.  Direct differentiation
of \eqref{eq:W} at the normalized umbilic gives the identity
\[
                              W_{10}=2q_{20}a^2(1-m),
\]
which vanishes when \(m=1\).  This completes the proof.
\end{proof}

\begin{corollary}[Jet-ideal form]                                  \label{cor:ideal}
Fix \(m\ge1\), and set
\begin{align*}
 \mathcal R_m={}&
 \R[p_{ij},q_{ij}:i+j\le m+2,p_{11}^{-1}]\big/
 \bigl(p_{01},q_{10},p_{10}-q_{01},p_{20},p_{02},q_{11},
       q_{02}+mp_{11}\bigr).
\end{align*}
This is the universal normalized jet ring, localized at
\(a=p_{11}\).  If
\[
             I_{<m}=\langle W_{ij}:i+j<m\rangle\subset\mathcal R_m,
\]
then
\[
                              W_{m0}\in I_{<m}.
\]
\end{corollary}

\begin{proof}
For \(m=1\), the identity \(W_{10}=2q_{20}a^2(1-m)\) proves the claim
directly.  Let \(m>1\), and carry out the preceding finite normal-form
calculation in the quotient by
\(I_{<m}\).  The hodograph coefficients require division only by powers of
the unit \(a\) (equivalently, by \(J(0)=(m+1)a^2/2\)).  Every homological
elimination divides only by one of the nonzero numerical eigenvalues in
\eqref{eq:homA}--\eqref{eq:homB}.  In the odd case, equation
\eqref{eq:killc} sets the sole resonant coefficient equal to zero in that
quotient.  Hence the coefficient in \eqref{eq:return} is zero in the
quotient, which is exactly the asserted ideal membership.
\end{proof}

\section{Consequences and status}

For an analytic \(W\)-congruence the literal equation \(W_{m0}=0\) is, of
course, immediate from the definition \(W\equiv0\).  Theorem~\ref{thm:main}
proves the stronger assertion intended in the jet-space discussion of
\cite{CG}: before one assumes the order-\(m\) equation, it is already forced
by the equations of smaller total order.  Corollary~\ref{cor:ideal} gives
the corresponding algebraic redundancy.  Thus the exceptional Taylor
equation introduces no additional jet relation for any integer \(m\ge1\).

The proof also explains the low-order computer algebra observed in the
source.  The resonance is not removed by the coefficient that would normally
serve as a pivot.  Instead, the exact identity \eqref{eq:key} converts a
lower mixed equation into the compatibility condition that removes the only
possible nonlinear normal-form obstruction.

Accordingly, the conjecture represented by equation~(3.4) in the current
version of \cite{CG} is true, both in its literal form and in the stronger
finite-jet redundancy sense made precise in Theorem~\ref{thm:main}.

\begin{thebibliography}{9}
\bibitem{CG}
M.~Craizer and R.~A.~Garcia,
\emph{Stable umbilical points of a \(W\)-congruence},
arXiv:2504.07289v3 (26 December 2025),
\href{https://arxiv.org/abs/2504.07289}{arXiv:2504.07289}.
\end{thebibliography}

\end{document}
